\section{Conclusion}
\label{sec:conclusion}
This article covered the full path from controller design to embedded
validation for the fixed-wing flight mode of a hybrid VTOL aircraft. A
cascade PID and an output-feedback LQR were designed on the linearized
model, verified on the nonlinear model, stressed under model mismatch
and wind, and deployed to an STM32 microcontroller through automatic
code generation, so that the artifact exercised on the bench is the
production code itself.

Four results stand out. First, on the nominal plant the
requirements-retuned PID dominates the performance comparison: zero
overshoot and the tighter tracking on every guidance channel, meeting
every requirement row, while the deliberately soft fixed-gain LQRy
concedes a factor of ${\approx}5$ in airspeed RMS and violates the
airspeed-overshoot requirement ($16.5\%$). Second, the robustness
study inverts that verdict where it matters most: under a severe,
sign-changing lateral--directional model mismatch the PID loses the
mission outright, while the fixed-gain LQRy completes it with
near-nominal error --- a tolerance consistent with the classical
robustness heritage of LQ designs, earned here by the deliberately
soft tuning; the remaining scenarios add nuance: under the
longitudinal-only mismatch
the LQRy stays nominal while the PID's clamped altitude loop drifts,
under wind and gusts the PID holds every channel more tightly, and
both tolerate a plausible $+10\%$ identification uncertainty. The
comparison thus isolates a genuine performance--robustness trade-off
between the two standard laws: which one is preferable depends on
whether the dominant threat is model error --- where the PID's
failure is catastrophic and the LQRy's tolerance is decisive --- or
atmospheric disturbance, where the PID holds the line more tightly.
Third, the lockstep HIL architecture
is practically transparent when the campaign is scripted from a single
source of references and the controller start-up is synchronized: for
the synchronized PID deployment, the SIL--HIL residual reaches the
quantization floor of the communication
bus ($0.011^\circ$ in pitch attitude in the worst flight case), and the
loop is bit-deterministic --- the same generated code executed on two
different microcontroller models reproduces full 300-second flight
cases bit for bit --- which makes the bench a measuring instrument and allows a
campaign to be reproduced exactly on different hardware. Fourth,
exploiting that property, a sensitivity study ranked the implementation
choices that actually matter for SIL--HIL agreement: reference
consistency between campaigns dominates by two orders of magnitude,
start-up desynchronization leaves a persistent residual that link
quality cannot remove, low-rate frame loss is benign, and controller
discretization --- the usual first suspect --- is negligible at a
10\,ms sample rate.

Future work follows directly from the limitations identified here: a
feedforward decoupling term for the roll-to-pitch coupling quantified
in Section~\ref{sec:coupling}, evaluated on both control laws; an
extended robustness campaign with stochastic Dryden/von K\'arm\'an
turbulence~\cite{milf8785c1980,milhdbk1797} in addition to the
discrete gusts used here, and with robust syntheses designed for this
trade-off~\cite{abdulhamid2016robusto} as a third contender;
equalization of the two design harnesses (actuator models, reference
shaping, and protection logic) so that future comparisons are free of
these asymmetries; a requirements set rederived specifically for
unmanned aircraft in the sense of the scaled flying-qualities
literature~\cite{cotting2010evolution,greene2014toward};
extension of the LQRy with gain scheduling over the flight envelope,
evaluated under this same SIL-to-HIL protocol; and flight
testing of the embedded controllers validated on the bench, closing
the simulation-to-flight loop that the HIL stage
approximates~\cite{shen2025model}.
